% Generated draft: requires mathematical and source review.
\documentclass[11pt]{article}
\usepackage[margin=25mm]{geometry}
\usepackage{fontspec}
\setmainfont{Latin Modern Roman}
\usepackage{amsmath,amssymb,mathtools,unicode-math}
\setmathfont{Latin Modern Math}
\usepackage{enumitem,needspace,xurl,hyperref}
\setlength{\parindent}{0pt}
\setlength{\parskip}{5pt}
\setlength{\emergencystretch}{3em}
\newcommand{\N}{\mathbb{N}}
\newcommand{\Z}{\mathbb{Z}}
\newcommand{\Q}{\mathbb{Q}}
\newcommand{\R}{\mathbb{R}}
\newcommand{\C}{\mathbb{C}}
\newcommand{\F}{\mathbb{F}}
\newcommand{\A}{\mathbb{A}}
\newcommand{\PP}{\mathbb{P}}
\newcommand{\E}{\mathbb{E}}
\newcommand{\eps}{\varepsilon}
\newcommand{\dd}{\,\mathrm{d}}
\begin{document}
\section*{UnsolvedMath ID: 801033}
\textbf{Source problem code:} catalogue entry
\subsection{Definitions and mathematical statement}
A computable metric presentation is a countable metric space \(M=(\{m_i:i\in\N\},d_M)\) whose distances \(d_M(m_i,m_j)\) are uniformly computable real numbers. Write \(\overline M\) for its metric completion. Points of \(\overline M\) are represented by effective Cauchy names: sequences \((i_n)_{n\in\N}\) satisfying \(d_M(x,m_{i_n})\leq 2^{-n}\). A Polish space is a separable completely metrizable topological space; here a computable Polish space means one admitting a computable metric presentation. The symbol \(\cong\) in the supplied question is interpreted as homeomorphism, as indicated by its topological formulation.

For a Turing degree \(\alpha\), a map between completions is \(\alpha\)-computable if an algorithm using an oracle of degree \(\alpha\) transforms every name of an input point into a name of its image. The degree \(\mathbf 0\) is the computable Turing degree. The source calls \(X\) \(\alpha\)-categorical if every pair of computable metric presentations whose completions are homeomorphic to \(X\) admits an \(\alpha\)-computable homeomorphism between the completions. If a least such degree exists in the Turing-degree ordering, the source calls it the degree of topological categoricity of \(X\); existence of a least degree is not assumed here.

\textbf{Question.} Does there exist an infinite computable Polish space \(X\) such that
\[
\forall M,N\;\Bigl[
\bigl(M,N\text{ are computable metric presentations}\bigr)
\land \overline M\cong X\land\overline N\cong X
\ \Longrightarrow\
\exists h:\overline M\longrightarrow\overline N\;
\bigl(h\text{ is a computable homeomorphism}\bigr)
\Bigr]?
\]
The quantification over presentations does not require a computable witness to either homeomorphism with \(X\), and the question does not request a uniform procedure producing \(h\) from presentations.

\textbf{Unresolved convention.} The supplied record does not specify whether a ``computable homeomorphism'' requires computability of both \(h\) and \(h^{-1}\), or only of the forward map \(h\). Under the two-sided interpretation, replace the final predicate above by
\[
 h\text{ is a homeomorphism},\qquad
 h\text{ is computable},\qquad h^{-1}\text{ is computable}.
\]
Under the forward-only interpretation, omit the last condition. The intended convention requires confirmation from the source before a unique fully specified formulation can be accepted. The effective Cauchy-name framework above supplies a standard interpretation of computability that is not explicitly detailed in the supplied record.
\subsection{Short English statement}
Does some infinite computably presented Polish space have the property that every two computable presentations of its topology are computably homeomorphic? The source labels this question open. Whether the homeomorphism must also have a computable inverse remains unspecified in the supplied record.
\subsection{Sources}
\begin{itemize}
\item Supplied local UnsolvedMath registry record: permanent ID \(801033\), exact source code \texttt{catalogue entry}, title \emph{Open Problems in Mathematical Logic --- Problem 45}. The registry attributes the problem to Barmpalias, Gao et al., with proposed year 2026, and labels it \texttt{open}; this is a supplied status, not an independently verified assessment.
\item Supplied source attribution: Barmpalias, Gao et al., \emph{Open Problems in Mathematical Logic} (2026), Problem 45, \url{https://arxiv.org/abs/2608.26628}. The background records access on 2026-10-16 and extraction by \texttt{arxiv\_latex}. These are supplied provenance claims; the source was not independently accessed or verified for this draft.
\item Supplied upstream revision: \texttt{d371ee917ac9647fbb06cd6fc529deec6f3dda82}. The registry records its update time as \texttt{2026-10-16T14:00:00Z}.
\end{itemize}
\end{document}
